\chapter{Mô tả tổng quan hệ thống}

\section{Giới thiệu}

AITasker là nền tảng trực tuyến hỗ trợ kết nối giữa doanh nghiệp có nhu cầu
triển khai các giải pháp trí tuệ nhân tạo và các chuyên gia AI có chuyên môn
phù hợp để thực hiện dự án.

Hệ thống hỗ trợ toàn bộ quy trình hợp tác, bao gồm đăng ký tài khoản, đăng nhập,
tạo dự án, tìm kiếm dự án, gửi đề xuất, chấp nhận đề xuất, tạo hợp đồng,
quản lý mốc công việc, thanh toán, trao đổi tin nhắn, gửi thông báo và đánh giá
sau khi hợp đồng hoàn thành.

AITasker được xây dựng dưới dạng ứng dụng Web. Frontend sử dụng ReactJS,
Backend sử dụng FastAPI, dữ liệu được lưu trữ trong PostgreSQL và được ánh xạ
thông qua SQLAlchemy ORM. Hệ thống sử dụng JWT để xác thực và BCrypt để băm
mật khẩu người dùng.

%==============================================================

\section{Mục tiêu của hệ thống}

Hệ thống AITasker được phát triển nhằm đạt được các mục tiêu sau:

\begin{enumerate}
    \item Kết nối doanh nghiệp với các chuyên gia AI có năng lực phù hợp.
    \item Chuẩn hóa quy trình tuyển chọn và thuê chuyên gia AI.
    \item Hỗ trợ doanh nghiệp tạo và quản lý dự án trên một nền tảng thống nhất.
    \item Hỗ trợ chuyên gia tìm kiếm dự án và gửi đề xuất thực hiện.
    \item Quản lý hợp đồng, mốc công việc và tiến độ triển khai dự án.
    \item Minh bạch hóa quá trình thanh toán và lịch sử giao dịch.
    \item Hỗ trợ trao đổi thông tin giữa doanh nghiệp và chuyên gia.
    \item Cung cấp cơ chế đánh giá sau khi hoàn thành hợp đồng.
    \item Hỗ trợ quản trị viên giám sát và quản lý hoạt động toàn hệ thống.
\end{enumerate}

%==============================================================

\section{Kiến trúc tổng quan hệ thống}

AITasker được thiết kế theo kiến trúc phân lớp, bao gồm lớp giao diện,
lớp xử lý nghiệp vụ và lớp dữ liệu.

\begin{figure}[H]
    \centering
    \includegraphics[
        width=\textwidth,
        height=0.78\textheight,
        keepaspectratio
    ]{architecture.png}
    \caption{Kiến trúc tổng quan hệ thống AITasker}
    \label{fig:architecture-overview}
\end{figure}

Như thể hiện tại Hình~\ref{fig:architecture-overview}, ba nhóm người dùng gồm
doanh nghiệp, chuyên gia AI và quản trị viên truy cập hệ thống thông qua ứng
dụng Web ReactJS.

Frontend gửi các yêu cầu HTTP đến REST API được xây dựng bằng FastAPI.
Backend tiếp nhận yêu cầu, thực hiện xác thực, kiểm tra quyền truy cập và điều
phối yêu cầu đến các dịch vụ nghiệp vụ tương ứng.

Các dịch vụ nghiệp vụ truy cập PostgreSQL để lưu trữ và truy vấn dữ liệu.
Hệ thống có thể được triển khai trên môi trường điện toán đám mây để tăng
khả năng sẵn sàng và mở rộng.

%--------------------------------------------------------------

\subsection{Lớp giao diện người dùng}

Lớp giao diện được phát triển bằng ReactJS và cung cấp các màn hình phù hợp
với vai trò của từng người dùng.

Các chức năng giao diện chính gồm:

\begin{itemize}
    \item Đăng ký và đăng nhập.
    \item Quản lý hồ sơ cá nhân.
    \item Quản lý dự án.
    \item Tìm kiếm dự án và chuyên gia.
    \item Quản lý đề xuất.
    \item Quản lý hợp đồng và mốc công việc.
    \item Quản lý thanh toán.
    \item Trao đổi tin nhắn.
    \item Xem thông báo.
    \item Đánh giá người dùng.
    \item Xem Dashboard và báo cáo thống kê.
\end{itemize}

%--------------------------------------------------------------

\subsection{Lớp xử lý nghiệp vụ}

Backend được xây dựng bằng FastAPI và được chia thành các dịch vụ nghiệp vụ.

Các thành phần chính gồm:

\begin{itemize}
    \item Authentication Service.
    \item User Service.
    \item Enterprise Service.
    \item Expert Service.
    \item Category Service.
    \item Skill Service.
    \item Project Service.
    \item Proposal Service.
    \item Contract Service.
    \item Milestone Service.
    \item Payment Service.
    \item Conversation Service.
    \item Message Service.
    \item Notification Service.
    \item Review Service.
    \item Dashboard Service.
    \item Statistics Service.
\end{itemize}

%--------------------------------------------------------------

\subsection{Lớp dữ liệu}

Hệ thống sử dụng PostgreSQL làm hệ quản trị cơ sở dữ liệu và SQLAlchemy làm
công cụ ánh xạ đối tượng quan hệ.

Các nhóm dữ liệu chính gồm:

\begin{itemize}
    \item Người dùng, vai trò và quyền.
    \item Hồ sơ doanh nghiệp và chuyên gia.
    \item Danh mục và kỹ năng.
    \item Dự án và tệp đính kèm.
    \item Đề xuất và hợp đồng.
    \item Mốc công việc và thanh toán.
    \item Hội thoại và tin nhắn.
    \item Thông báo và đánh giá.
\end{itemize}

%==============================================================

\section{Các tác nhân của hệ thống}

Hệ thống có ba tác nhân chính:

\begin{itemize}
    \item Quản trị viên.
    \item Doanh nghiệp.
    \item Chuyên gia AI.
\end{itemize}

Biểu đồ Use Case tổng thể được trình bày trong
Hình~\ref{fig:usecase-overview}.

\begin{figure}[H]
    \centering
    \includegraphics[
        width=0.82\textwidth,
        height=0.85\textheight,
        keepaspectratio
    ]{Usecase.png}
    \caption{Biểu đồ Use Case tổng thể của hệ thống AITasker}
    \label{fig:usecase-overview}
\end{figure}

%--------------------------------------------------------------

\subsection{Quản trị viên}

Quản trị viên chịu trách nhiệm quản lý và giám sát hệ thống.

Các chức năng chính gồm:

\begin{itemize}
    \item Quản lý người dùng.
    \item Quản lý danh mục.
    \item Quản lý các dự án trên hệ thống.
    \item Xem báo cáo và số liệu thống kê.
    \item Theo dõi tình trạng hoạt động của hệ thống.
\end{itemize}

%--------------------------------------------------------------

\subsection{Doanh nghiệp}

Doanh nghiệp là tác nhân tạo dự án và thuê chuyên gia AI.

Các chức năng chính gồm:

\begin{itemize}
    \item Đăng ký và đăng nhập.
    \item Quản lý hồ sơ cá nhân.
    \item Tạo, cập nhật và xóa dự án.
    \item Xem danh sách và chi tiết dự án.
    \item Xem, chấp nhận hoặc từ chối đề xuất.
    \item Tạo và quản lý hợp đồng.
    \item Tạo mốc công việc.
    \item Theo dõi tiến độ.
    \item Thanh toán theo mốc công việc.
    \item Gửi tin nhắn và xem hội thoại.
    \item Nhận thông báo.
    \item Đánh giá chuyên gia.
\end{itemize}

%--------------------------------------------------------------

\subsection{Chuyên gia AI}

Chuyên gia AI là tác nhân tìm kiếm và thực hiện các dự án do doanh nghiệp đăng.

Các chức năng chính gồm:

\begin{itemize}
    \item Đăng ký và đăng nhập.
    \item Quản lý hồ sơ và kỹ năng.
    \item Tìm kiếm và xem chi tiết dự án.
    \item Gửi hoặc rút đề xuất.
    \item Xem và quản lý hợp đồng.
    \item Cập nhật tiến độ thực hiện.
    \item Gửi tin nhắn và xem hội thoại.
    \item Xem lịch sử thanh toán.
    \item Nhận thông báo.
    \item Đánh giá doanh nghiệp.
\end{itemize}

%==============================================================

\section{Quy trình nghiệp vụ tổng quát}

Quy trình hợp tác chính của hệ thống được thực hiện như sau:

\begin{enumerate}
    \item Người dùng đăng ký và đăng nhập vào hệ thống.
    \item Doanh nghiệp tạo một dự án mới.
    \item Chuyên gia tìm kiếm và xem thông tin dự án.
    \item Chuyên gia gửi đề xuất cho dự án.
    \item Doanh nghiệp xem và lựa chọn một đề xuất phù hợp.
    \item Hệ thống chấp nhận đề xuất và tạo hợp đồng.
    \item Doanh nghiệp và chuyên gia thực hiện công việc theo các Milestone.
    \item Doanh nghiệp thanh toán sau khi Milestone hoàn thành.
    \item Khi toàn bộ công việc hoàn thành, hợp đồng chuyển sang trạng thái
    \texttt{COMPLETED}.
    \item Hai bên thực hiện đánh giá lẫn nhau.
\end{enumerate}

%==============================================================

\section{Môi trường hoạt động}

Hệ thống hoạt động trên nền Web và hỗ trợ các thành phần sau:

\begin{itemize}
    \item Hệ điều hành Windows, Linux hoặc macOS.
    \item Trình duyệt Google Chrome, Microsoft Edge và Mozilla Firefox.
    \item Python 3.12.
    \item FastAPI.
    \item ReactJS.
    \item PostgreSQL.
    \item SQLAlchemy.
    \item JWT và BCrypt.
\end{itemize}

%==============================================================

\section{Các ràng buộc nghiệp vụ}

Các ràng buộc chính của hệ thống gồm:

\begin{itemize}
    \item Người dùng phải đăng nhập để sử dụng các chức năng nghiệp vụ.
    \item Email và tên đăng nhập phải là duy nhất.
    \item Chỉ doanh nghiệp được phép tạo dự án.
    \item Chỉ chuyên gia được phép gửi đề xuất.
    \item Chuyên gia không được gửi nhiều đề xuất cho cùng một dự án.
    \item Đề xuất chỉ được gửi khi dự án đang ở trạng thái \texttt{OPEN}.
    \item Khi một đề xuất được chấp nhận, các đề xuất còn lại phải bị từ chối.
    \item Chấp nhận đề xuất phải dẫn đến việc tạo hợp đồng.
    \item Chỉ hợp đồng \texttt{ACTIVE} mới được tiếp tục thực hiện.
    \item Milestone phải hoàn thành trước khi được thanh toán.
    \item Chỉ hợp đồng \texttt{COMPLETED} mới được đánh giá.
    \item Mỗi người dùng chỉ được đánh giá một lần cho cùng một hợp đồng.
    \item Mật khẩu phải được băm trước khi lưu vào cơ sở dữ liệu.
    \item Access Token phải có thời hạn sử dụng.
\end{itemize}

%==============================================================

\section{Giả định và phụ thuộc}

Hệ thống được xây dựng dựa trên các giả định sau:

\begin{itemize}
    \item Người dùng có kết nối Internet ổn định.
    \item Frontend có thể kết nối đến REST API của Backend.
    \item Máy chủ PostgreSQL hoạt động bình thường.
    \item Dịch vụ lưu trữ tệp luôn khả dụng.
    \item Hệ thống triển khai thực tế sử dụng kết nối HTTPS.
    \item Người dùng cung cấp thông tin chính xác khi tạo hồ sơ và dự án.
\end{itemize}